\chapter{Development Workflow}
\label{chap-dev-workflow}

\renewcommand{\headrulewidth}{0.3pt}
\fancyhead[LE]{\thepage} \fancyhead[RE]{\nouppercase{\textbf{\leftmark}}}
\fancyhead[RO]{\thepage} \fancyhead[LO]{\nouppercase{\textbf{\rightmark}}}
\fancyfoot[L,C,R]{}

\setlength{\parindent}{0.35cm}

\minitoc
% ================================================================

This chapter gathers the \textbf{rules and processes} a developer must follow when working on \cwv: how to extend the application with a new operation, the checklist to run through before committing, the project conventions, and how versioning and code documentation are handled.

\section{Adding a New Dialog Operation}
\label{sec-dev-add-op}

To add a new dialog-level workflow, follow three steps:

\begin{enumerate}
    \item define a \script{*Result} dataclass in \script{ht/client/api\_types.py};
    \item implement \script{run\_<name>} in \script{ht/client/operations\_client.py} --- it owns the \script{fetch}~$\rightarrow$~\script{transform}~$\rightarrow$~\script{render} orchestration;
    \item add a one-line facade method on \script{HydrologicalTwinClient}, then call it from \script{ExploreData}.
\end{enumerate}

Method bodies on the client class are intentionally thin ($\leq$~5 statements). Keep them that way.

\section{When Editing --- Checklist}
\label{sec-dev-checklist}

\begin{itemize}
    \item touching dialog code? $\rightarrow$ \script{frontend/Dialog*.py}. Route through \script{ExploreData};
    \item adding a new dialog-level workflow? $\rightarrow$ the three-step recipe above (\cref{sec-dev-add-op});
    \item adding a new compute primitive? $\rightarrow$ \script{ht/developer/handlers.py} (compute) or \script{services/public/twin\_io.py} (L3 state reads / disk I/O); route \textit{via} \script{dispatch.py} with a new \script{kind=};
    \item found a \script{qgis} / \script{PyQt5} import in \script{external/HydrologicalTwinAlphaSeries/}? $\rightarrow$ bug. Move it to \script{interface/} or \script{frontend/};
    \item found a module under \script{services/} importing from \script{ht/client/} or \script{ht/developer/}? $\rightarrow$ bug. That is an \emph{upward} import; imports must point downward only;
    \item found a frontend module reaching past \script{ExploreData} into HTAS internals? $\rightarrow$ bug. Route through the client gate;
    \item found data handling inside \script{ht/client/} (\script{DataFrame} building, array reshaping/aggregation, unit/path shaping, disk I/O, domain math)? $\rightarrow$ bug. L1 only orchestrates.
\end{itemize}

\section{Conventions}
\label{sec-dev-conventions}

\begin{itemize}
    \item \textbf{Environment:} \script{pixi}-managed (\script{pixi.toml}, \script{pixi.lock}). QGIS, Python, and all dependencies are pinned in a single reproducible environment (\cref{sec-install-pixi});
    \item \textbf{Sync helpers:} \script{./verifState.sh} (diagnostics) and \script{./updateProject.sh} (ordered multi-repository / submodule sync);
    \item \textbf{Pre-lifecycle discovery} (no twin required): \script{HydrologicalTwinClient.detect\_from\_out\_caw(...)} and \script{detect\_project\_neighbors(...)} --- static helpers used by \script{DialogSettings} to auto-fill paths before \script{configure} / \script{load};
    \item \textbf{Compatibility shim:} \script{cawaqsviz/backend}, if present, re-exports HTAS symbols. The real implementation lives only in \script{external/HydrologicalTwinAlphaSeries/}.
\end{itemize}

\section{Version Update}

The plugin version is specified in the \script{metadata.txt} file contained in code excerpt \ref{code_metadata}.
\begin{lstlisting}[language=json, caption={Mandatory information to fill in the \script{metadata.txt} file}, label={code_metadata}]
[general]
name=CaWaQS-Viz
qgisMinimumVersion=3.0
description=Description
version=0.1.16
author=Lise-Marie GIROD (Mines Paris - PSL), Nicolas Flipo (Mines Paris - PSL)
email=lise-marie.girod@minesparis.psl.eu

about=Python QGIS-based vizualization software for CaWaQS3.x

repository=https://gitlab.com/gviz/cawaqsviz
\end{lstlisting}

\section{Code Documentation}

The code documentation is published on the following link via GitLab Pages: \href{https://cawaqsviz-gviz-adf347ff6d9d45c174db091497a2dfc0f14305cdb4ed9070.gitlab.io/}{https://cawaqsviz-gviz-adf347ff6d9d45c174db091497a2dfc0f14305cdb4ed9070.gitlab.io/}. The documentation is maintained through docstrings in the application modules, following the formatting shown in the example below.


\begin{lstlisting}[language=Python, caption={Docstring formatting for automatic integration into code documentation}, label={code_format_docstrings}]
class Exemple:
    """
    Resume des fonctionnalites de la classe.
    """

    def __init__(self, arg1):
        """
        Methode constructeur.

        :param arg1: descriptif du parametre
        :type arg1: str
        """
        self.arg1 = arg1

    def do_something(self, arg1: list) -> str:
        """
        Descriptif de la fonction.

        :param arg1: descriptif du parametre
        :type arg1: list
        :return: descriptif de ce que la fonction retourne
        :rtype: str
        """
        return ", ".join(map(str, arg1))
\end{lstlisting}
\vspace{1em}
Compilation is performed by \script{Sphinx}\footnote{whose documentation is accessible via \href{https://www.sphinx-doc.org/en/master/}{https://www.sphinx-doc.org/en/master/}}. The \script{docs/sources} directory contains the configuration files (\script{Conf.py}), as well as the documentation content.
When integrating new modules into the plugin, adding them to the index (\script{index.rst}) is necessary for the code docstring to appear on the documentation page. This index integrates 3 files: \script{frontend.rst}, \script{interface.rst}, \script{backend.rst}, which contain the classes to be documented. Integrating a new module into the application requires it to be added to these files in the format of code excerpt \ref{code_format_rst}. Updates to the documentation source files must be committed and synchronised with the repositories.

\begin{lstlisting}[language=Python,caption={Docstring formatting for automatic integration into code documentation}, label={code_format_rst}]
Module
------

.. automodule:: cawaqsviz.Module
   :members:
   :show-inheritance:
   :undoc-members:
\end{lstlisting}

\vspace{1em}

Finally, with each new commit to the \script{Main} or
\script{Beta} branches, the documentation is compiled according to the content of the \script{docs/source} directory of the repository.

\begin{lstlisting}
image: python:3.9

stages:
  - build
  - deploy

before_script:
  - pip install -U pip
  - pip install -r requirements.txt

build-docs:
  stage: build
  script:
    - sphinx-build -b html docs/source public
  artifacts:
    paths:
      - public

pages:
  stage: deploy
  script:
    - echo "Pages deployed"
  artifacts:
    paths:
      - public
  only:
    - main
    - beta

\end{lstlisting}
